\section{模型总结与迁移}
\ifteachernotes
\block{一 先判断公共点的身份}

这里把等腰三角形两腰所夹的顶角叫作“手”，底边两端的底角顶点叫作“脚”。在本章的等腰直角三角形中，“手”就是直角顶点。

{\small\renewcommand{\arraystretch}{1.25}
\begin{tabularx}{\linewidth}{@{}p{1.2cm} X X p{2cm}@{}}
\toprule
构型 & 公共点的身份 & 优先考虑的构造 & 本章对应 \\
\midrule
手拉手 & 两个三角形的直角顶点 & 绕公共顶点同向旋转 \(90^\circ\)，配对端点 & 两道题构造后的新图形 \\
手拉脚 & 一个直角顶点、一个底角顶点 & 倍长经过直角顶点的腰，或旋转另一条边 & 题 1 \\
脚拉脚 & 两个底角顶点 & 各倍长一条腰，分别把公共点补成直角顶点 & 题 2 \\
\bottomrule
\end{tabularx}}

名称只是帮助认图。判断时要写出实际的等长条件、直角位置和旋转对应，不能仅凭图形像不像。

\block{二 手拉手为什么能同时得到等长与垂直}

设 \(\triangle OAB\)、\(\triangle OCD\) 是共用直角顶点 \(O\) 的等腰直角三角形，\(OA=OB\)，\(OC=OD\)。若同一个方向的 \(90^\circ\)旋转使 \(A\) 对应 \(B\)，同时使 \(C\) 对应 \(D\)，那么整条 \(AC\) 就会对应 \(BD\)。

当 \(A\)、\(C\) 不重合时，旋转保持长度，又把线段方向转过 \(90^\circ\)，所以

\[
AC=BD,\qquad AC\perp BD.
\]

用全等来写，就是 \(\triangle OAC\cong \triangle OBD\)（\(\mathrm{SAS}\)）；用旋转来理解，还能立即说明垂直性。这也解释了为什么写证明时必须明确端点的对应。

题 1 方法二中，以 \(A\) 为中心，\(E\to D\)、\(C\to G\)，得到 \(CE=DG\)、\(CE\perp DG\)；题 2 中，以 \(C\) 为中心，\(B\to G\)、\(E\to F\)，得到 \(BE=GF\)，并得到对应角相等。随后还要利用中点，把这个新结论送到题目要求的线段上。

\block{三 倍长为什么能把脚补成手}

设 \(\triangle OXY\) 在 \(X\) 处为直角，\(XO=XY\)。公共点 \(O\) 原来是一个底角顶点。延长 \(YX\) 到 \(Z\)，使 \(XZ=XY\)，连接 \(OZ\)。

由于 \(XO=XY=XZ\)、\(XO\perp YZ\)，两侧小三角形都是等腰直角三角形，故

\[
OY=OZ,\qquad
\angle YOZ=45^\circ+45^\circ=90^\circ.
\]

新三角形 \(\triangle YOZ\) 的直角顶点便是 \(O\)。操作时应倍长经过原直角顶点 \(X\) 的腰，连接原底角顶点 \(O\) 与新的端点 \(Z\)。

题 1 方法一用 \(\triangle ABC\)、\(\triangle GBC\) 在 \(C\) 处拼出直角；方法二用 \(\triangle ABC\)、\(\triangle ABG\) 在 \(A\) 处拼出直角。题 2 则分别在 \(A\)、\(D\) 两个位置倍长，最后都在 \(C\) 处拼出直角。

\block{四 中点怎样帮助完成第二次搬运}

\textbf{两条线段共用中点。} 若 \(B\) 同时是 \(DF\)、\(AG\) 的中点，绕 \(B\) 半转便有 \(D\to F\)、\(A\to G\)，得到对应线段相等且平行。题 1 方法一采用的就是这个构造。

\textbf{中点配平行线。} 若 \(D\) 是 \(EF\) 的中点，\(G\) 关于 \(D\) 的对称点是 \(P\)，则半转有 \(F\to E\)、\(G\to P\)，得到 \(GF=EP\)、\(GF\parallel EP\)。题 2 把旋转型全等给出的 \(BE=GF\) 接到这里，才构造出 \(BE=EP\) 的等腰三角形。

遇到中点，不必机械地延长待证线段。先确定需要搬运哪条边，再决定应当倍长或作平行线的位置。

\block{五 与夹半角的联系}

题 1 的 \(A\) 处既有 \(\angle DAE=90^\circ\)，又有 \(\angle BAC=45^\circ\)。这提示可以把相关三角形绕 \(A\) 旋转 \(90^\circ\)，把原本分开的边放在一起。旋转之后，等腰直角三角形保证端点对应，中点条件完成另一组全等。

两类模型都要看等长和角度，但本章的结论来自具体的旋转对应；每一次换构型，都应重新检查对应端点和射线方向。

\block{易错提醒}

\begin{itemize}
\item 共用一个点，不等于共用直角顶点。先标明公共点在两个三角形中分别是什么身份。
\item 两个等腰直角三角形必须按同一旋转方向配对端点。把其中一个三角形翻到另一侧后，原来的连线结论可能改变。
\item 全等判定要根据已有条件填写。本章关键步骤使用两边及夹角，判定是 \(\mathrm{SAS}\)；平行线构造中的两角及夹边可以用 \(\mathrm{ASA}\)。
\item 中点带来的等长、平行必须由全等或半转说明，不能仅因图中出现“八字形”就直接使用。
\item 等长与垂直要分别交代依据。只证明了 \(CE=CF\)，还不能据此断定 \(CE\perp CF\)。
\item 题 2 的“如图”包含两个三角形的方位。本章将它明确为 \(E\)、\(C\) 在 \(AB\) 两侧；若把 \(E\) 换到另一侧，\(BE\perp BC\) 不再是同一个命题。
\item 点序改变后，重新判断角的加减；辅助三角形退化时，改用仍然有效的旋转或对称关系。
\end{itemize}
\else
请从公共点的身份、同向旋转、倍长与中点搬运四个方面，自主整理本章方法。
\fi
